\section{Model and Definitions}
\label{sec:model}

This section fixes the finite-state model used by both implementations and by
all exact-coverage measurements.  The definitions deliberately separate the
cryptographic primitive from the reduction and from the TMTO table structure.

\subsection{Reduced effective DES state space}

For an integer
\[
    1\leq b\leq 28,
\]
let
\[
    X_b=\{0,1,\ldots,2^b-1\},
    \qquad
    N=|X_b|=2^b.
\]
The parameter $b$ counts \emph{effective DES key bits}; it is not a count of
raw bits in the eight-byte DES key representation.

DES stores eight bytes, with one parity bit in each byte.  The experiment
keeps the first four bytes fixed and allows the final four bytes to carry at
most $4\cdot 7=28$ effective data bits.  Let
\[
    K_b:X_b\longrightarrow\{0,1\}^{64}
\]
denote the canonical encoding used by the artifact.  The state is decomposed
into four seven-bit groups, placed into the four variable DES bytes from most
significant group to least significant group, and the least-significant bit of
each byte is then chosen to give odd DES parity.  Unused high state bits are
zero.  This convention makes the state-to-key map deterministic and invertible
on its image.

The fixed four-byte prefix uses seven effective one-bits in each byte.  Its
purpose is solely to preserve the effective-data structure of the historical
prototype; it is not a security parameter.

\subsection{Encryption oracle and reduction}

Let $P$ be the fixed ASCII plaintext
\[
  \texttt{HPC reproducibility fixture}.
\]
For $x\in X_b$, define the deterministic encryption oracle
\[
    G_b(x)=
    \operatorname{DES}_{K_b(x)}(P),
\]
using DES in ECB mode with PKCS\#5 padding.  DES is obsolete and appears here
only as a small deterministic fixture whose entire reduced state space can be
enumerated.

The reduction
\[
    R_b:\{0,1\}^{*}\longrightarrow X_b
\]
uses the final four ciphertext bytes.  From each such byte it discards the
least-significant bit, retains the upper seven bits, concatenates the resulting
four seven-bit words into a 28-bit integer, and finally truncates to the low
$b$ bits.  Equivalently, if the final four ciphertext bytes are
$c_0,c_1,c_2,c_3$, interpreted as unsigned integers, define
\[
    u_i=\left\lfloor c_i/2\right\rfloor\in\{0,\ldots,127\},
\]
\[
    U =
       (((u_0\cdot 2^7+u_1)\cdot 2^7+u_2)\cdot 2^7+u_3),
\]
and
\[
    R_b(c)=U\bmod 2^b.
\]

The TMTO iteration map is therefore
\begin{equation}
    F_b = R_b\circ G_b:X_b\longrightarrow X_b.
    \label{eq:iteration-map}
\end{equation}
For fixed $b$ and $P$, this is one deterministic functional graph on $N$
vertices.

\subsection{Targets and exact recovery}

A trial chooses a target state
\[
    x^\star\in X_b
\]
and publishes only
\[
    y^\star = G_b(x^\star).
\]
An algorithm reports exact recovery only when it returns a state $\hat x$ such
that
\[
    G_b(\hat x)=y^\star.
\]
The implementation therefore verifies candidate states against the full target
ciphertext rather than treating an endpoint collision as success.

For the deliberately reduced state spaces studied here, exhaustive search is a
reference method: test $x=0,1,\ldots,N-1$ until the target ciphertext is
matched.  It has exact state-space coverage one and serves as a correctness
control rather than as a practical attack on DES.

\subsection{Hellman matrix and table}

Choose $m$ distinct start states
\[
    S=(s_1,\ldots,s_m)\subset X_b.
\]
For a chain length $t\geq 1$, define the Hellman matrix
\[
    H_{i,0}=s_i,
    \qquad
    H_{i,j+1}=F_b(H_{i,j}),
    \qquad
    0\leq j<t.
\]
Thus row $i$ contains $t+1$ states
\[
    H_{i,0},H_{i,1},\ldots,H_{i,t},
\]
but the table represents the first $t$ positions
\[
    H_{i,0},\ldots,H_{i,t-1}
\]
and stores the pair
\[
    (H_{i,0},H_{i,t}).
\]
This convention matches the coverage definition used by Ma and
Hong~\cite{MaHong2009} and matches the artifact's exact coverage analyzer.

The implementation indexes all starts by their final endpoint.  Multiple rows
with the same endpoint are retained rather than silently overwritten.

\subsection{Exact Hellman coverage}

The represented-state set of one Hellman matrix is
\begin{equation}
    \mathcal C_H
    =
    \{H_{i,j}:1\leq i\leq m,\;0\leq j<t\}.
    \label{eq:hellman-coverage-set}
\end{equation}
Its exact covered-state count and state-space coverage fraction are
\[
    C_H=|\mathcal C_H|,
    \qquad
    \rho_H=\frac{C_H}{N}.
\]
The normalized coverage rate used in the classical literature is instead
\[
    \operatorname{CR}_H=\frac{C_H}{mt}.
\]
The two normalizations answer different questions:
\[
    \rho_H
      =
      \frac{mt}{N}\operatorname{CR}_H.
\]
When $mt=N$, as in the initial $b=16,m=1024,t=64$ stress configuration, these
two numerical values happen to coincide.  They should not otherwise be
conflated.

The artifact computes $C_H$ exactly with a bit set over all $N$ reduced states;
this analysis is performed outside the timed preprocessing and lookup paths.

\subsection{Endpoint collisions and candidate verification}

Let
\[
    e_i=H_{i,t}
\]
be the endpoint of row $i$.  The number of distinct stored endpoints is
\[
    E_H=\left|\{e_1,\ldots,e_m\}\right|.
\]
Endpoint multiplicity is retained because several starts may produce the same
endpoint.

During online lookup, the target ciphertext is reduced and projected forward
to possible endpoint positions.  If a projected value hits the endpoint
index, all previously unchecked rows stored under that endpoint are regenerated
from their starts.  An endpoint hit that fails to reproduce $y^\star$ is a
false candidate event.  We therefore distinguish:

\begin{itemize}
  \item \emph{projected endpoint matches}: online projections that hit a stored
        endpoint key;
  \item \emph{candidate rows checked}: distinct stored rows regenerated after
        such hits;
  \item \emph{exact recovery}: regeneration finds a state whose full
        ciphertext equals $y^\star$.
\end{itemize}

These counters measure concrete lookup work and should not be replaced by a
success/failure bit alone.

\subsection{Distinguished-point tables}

Fix a difficulty parameter
\[
    1\leq d\leq b
\]
and define
\[
    D_d(x)=1
    \quad\Longleftrightarrow\quad
    x\bmod 2^d=0.
\]
Thus a fraction exactly $2^{-d}$ of the state space is distinguished.  Under
a uniform-state heuristic, a fresh state is distinguished with probability
\[
    p_d=2^{-d},
\]
so an unconstrained geometric stopping model has mean
\[
    1/p_d=2^d
\]
transitions.

For a maximum chain length $L$, a distinguished-point row starts at $s_i$ and
iterates $F_b$ until the first transition whose output is distinguished:
\[
    \tau_i
      =
      \min\{j\in\{1,\ldots,L\}:D_d(F_b^j(s_i))=1\}.
\]
If no such $j$ exists, the row is \emph{truncated} and is not stored.  A
successful row stores its start and distinguished endpoint
\[
    (s_i,F_b^{\tau_i}(s_i)).
\]

The exact DP coverage used by the artifact is
\[
    \mathcal C_{\rm DP}
      =
      \bigcup_{\tau_i\ {\rm exists}}
      \{F_b^j(s_i):0\leq j<\tau_i\}.
\]
In particular, truncated rows are excluded from this represented-state set.
This convention is made explicit because other implementations may count
partially generated truncated rows differently.

\subsection{Measured quantities}

For each parameter configuration and trial, the study records as available:
\[
\begin{array}{ll}
\text{state-space size} & N,\\
\text{exact coverage} & C/N,\\
\text{exact recovery} & \mathbf 1[\hat x\ {\rm recovered}],\\
\text{stored rows} & \text{number of accepted table rows},\\
\text{distinct endpoints} & E,\\
\text{truncated DP rows} & T_{\rm trunc},\\
\text{endpoint/projection hits} & Q_{\rm end},\\
\text{candidate rows regenerated} & Q_{\rm cand},\\
\text{offline time} & T_{\rm off},\\
\text{online time} & T_{\rm on}.
\end{array}
\]
Later parameter sweeps additionally record functional-graph statistics such as
row merges, self-intersections, and cycle/component information.  Those
quantities are reported only after the corresponding instrumentation is
committed.
